% Retained all-residue variant: relations-4-1-allres-cut15.
% Data flow: pkatrain/production_position.py, all_residue_tokens and functions.
% Readouts: pkanet/allres.py, outputs and pair_logits.
\begin{figure}[tbp]
    \centering
    \begin{tikzpicture}[
        box/.style={draw=black!50,rounded corners=2pt,align=center,
            inner sep=5pt,font=\footnotesize},
        flow/.style={-{Stealth[length=2mm]},thick,draw=black!65},
        read/.style={-{Stealth[length=1.5mm]},draw=blue!55!black,thin},
        residue/.style={circle,draw=blue!65!black,fill=blue!12,
            minimum size=6.5mm,inner sep=0pt,font=\small},
        query/.style={rectangle,rounded corners=1pt,draw=orange!75!black,
            fill=orange!15,minimum size=6.5mm,inner sep=0pt,font=\small},
        note/.style={align=center,font=\footnotesize}]
        \node[box,text width=13.5cm] (input) at (0,0)
            {\textbf{Backbone coordinates and amino-acid identities}\\
             Node features and local-frame edge geometry};
        \node[box,text width=6.2cm] (bound) at (-3.6,-1.3)
            {\textbf{Bound complex, $AB$}\\Within- and cross-chain edges};
        \node[box,text width=6.2cm] (free) at (3.6,-1.3)
            {\textbf{Separated partners, $F$}\\Cross-chain edges masked};
        \draw[flow] (input.south -| bound.north) -- (bound.north);
        \draw[flow] (input.south -| free.north) -- (free.north);
        \node[note] at (0,-2.05)
            {Same coordinates; the mask applies to all attention stages};

        % The enclosure shares parameters across state passes, not across layers.
        \draw[draw=blue!45!black,fill=blue!2,rounded corners=3pt]
            (-7,-2.5) rectangle (7,-11.5);
        \coordinate (boundpass) at (-6,-2.5);
        \coordinate (freepass) at (6,-2.5);
        \draw[flow] (bound.south -| boundpass) -- (boundpass);
        \draw[flow] (free.south -| freepass) -- (freepass);
        \node[note,font=\small\bfseries] at (0,-2.85)
            {Shared GQT parameters across bound and separated passes};

        \node[box,text width=3.4cm,fill=white] (resenc) at (-4.8,-4.35)
            {\textbf{Residue encoder}\\
             Inductive, two blocks\\20~\AA{}; 4 then 1 maps};
        \draw[flow] (-4.8,-3.25) -- (resenc.north);
        \node[note,text width=3.5cm] at (4.8,-4.35)
            {\textbf{Residues $h_j$}\\Keys and values};
        \foreach \i/\x in {1/-1.65,2/-0.55,3/0.55,4/1.65} {
            \node[residue] (r\i) at (\x,-4.35) {$h_{\i}$};
            \node[query] (q\i) at (\x,-6.25) {$q'_{\i}$};
        }
        \draw[flow] (resenc.east) -- (r1.west);
        \foreach \i/\j in {1/1,1/2,2/1,2/2,2/3,3/2,3/3,3/4,4/3,4/4}
            \draw[read] (r\i.south) -- (q\j.north);
        \node[note,text width=3.5cm] at (4.8,-5.4)
            {\textbf{Bipartite attention}\\Residues $\to$ queries\\15~\AA{} + FF};
        \node[box,text width=3.4cm,fill=white] (init) at (-4.8,-6.25)
            {\textbf{Query initialisation}\\$q_i^0=h_i+e(a_i)$};
        \draw[flow] (resenc.south) -- node[left,font=\footnotesize] {$h_i$} (init.north);
        \draw[flow] (init.east) -- (q1.west);
        \node[note,text width=3.5cm] at (4.8,-6.45)
            {\textbf{Updated queries $q_i^\prime$}\\One per residue};

        \node[box,fill=white,text width=8.6cm] (qupdate) at (0,-7.55)
            {\textbf{Query--query attention + FF} within 15~\AA\\
             Final residue tokens $z_i$};
        \foreach \i in {1,2,3,4}
            \draw[flow] (q\i.south) -- (qupdate.north -| q\i.south);

        % All heads branch from final query tokens, not from encoder residues h.
        \node[box,fill=white,text width=6cm] (maps) at (-3.6,-9)
            {\textbf{Three representations per residue}\\
             Side chain $z_i$; learned terminal maps\\
             $T_N(z_i)$ and $T_C(z_i)$};
        \node[box,fill=white,text width=6cm] (struct) at (3.6,-9.65)
            {\textbf{Structural heads from $z_i$}\\
             Burial from each final token\\
             Contacts from normalised token pairs\\
             $[\,\tilde z_i\odot\tilde z_j,\;|\tilde z_i-\tilde z_j|\,]$};
        \draw[flow] (qupdate.south -| maps.north) -- (maps.north);
        \draw[flow] (qupdate.south -| struct.north) -- (struct.north);
        \node[box,fill=white,text width=6cm] (heads) at (-3.6,-10.55)
            {\textbf{Shared scalar heads}\\
             Shift $8\tanh(w^\top t+b)$; titratability\\
             $t\in\{z_i,T_N(z_i),T_C(z_i)\}$};
        \draw[flow] (maps.south) -- (heads.north);
        \node[box,draw=blue!55!black,text width=13.5cm] (shift) at (0,-12.2)
            {Add the group reference to obtain $\mathrm{p}K_a^{AB}$ and $\mathrm{p}K_a^{F}$\\
             \textbf{Binding-induced shift}\quad
             $\Delta\mathrm{p}K_a=\mathrm{p}K_a^{AB}-\mathrm{p}K_a^{F}$};
        \draw[flow] (heads.south) -- (shift.north -| heads.south);
    \end{tikzpicture}
    \caption[Shared GQT encoder and bipartite query readout]{The retained all-residue GQT architecture. Bound and separated passes share all model parameters and coordinates; the separated pass masks cross-chain edges at every attention stage. Local operations are shared across residues and proteins, while each attention layer has its own parameters. The inductive residue encoder applies two attention blocks with four and one relation maps, respectively. Blue circles denote its outputs $h_i$. Each query starts from $h_i$ plus a learned amino-acid embedding $e(a_i)$. Bipartite cross-attention reads the encoded residues into the query set, followed by query--query attention to produce final tokens $z_i$ (orange squares show the intermediate queries). Arrows between node sets show schematic information flow. Each attention block includes a residual feed-forward update (FF). Learned affine maps produce N- and C-terminal representations from $z_i$ before shared shift and titratability heads. Burial and contact heads also read the final tokens, with contacts scored from symmetric features of normalised token pairs. The separated contact readout therefore combines independently encoded partners without their relative placement. Group-specific references convert environmental shifts to state $\mathrm{p}K_a$ values; their bound-minus-separated difference excludes conformational relaxation.}
    \label{fig:jaxka-model-overview}
\end{figure}
